% ============================================================
% CAPITOLUL 1 — INTRODUCERE
% ============================================================

\chapter{Introducere}
\label{ch:introducere}

\section{Context și motivație}
\label{sec:context}

Bolile infecțioase reprezintă una dintre cele mai vechi și persistente provocări ale
umanității. De la pandemia de gripă din 1918 la epidemia de SARS din 2003 și,
mai recent, pandemia de COVID-19 declanșată în 2019, evenimentele epidemice de
amploare au demonstrat în mod repetat capacitatea lor de a perturba profund structurile
sociale, economice și sanitare ale societăților moderne
\cite{andersonmay1991, hethcote2000}.

Înțelegerea și controlul dinamicii bolilor infecțioase nu constituie doar o problemă
medicală, ci reprezintă o provocare interdisciplinară care îmbină epidemiologia,
matematica aplicată, informatica și știința datelor. La baza acestei înțelegeri stau
\emph{modelele matematice compartimentale}, instrumente formale care permit
cuantificarea proceselor de transmitere, predicția evoluției epidemiilor și
evaluarea intervențiilor de sănătate publică.

Modelul \textbf{SIR} (Susceptibili-Infectați-Recuperați), introdus de Kermack și
McKendrick în lucrarea fondatoare din 1927 \cite{kermack1927}, rămâne și astăzi
fundamentul teoretic al epidemiologiei matematice. Simplitatea sa elegantă ---
trei ecuații diferențiale ordinare cuplate --- ascunde o bogăție matematică
remarcabilă și o putere predictivă validată de zeci de ani de aplicații empirice.

Odată cu creșterea exponențială a volumului de date epidemiologice disponibile și
cu diversificarea instrumentelor computaționale moderne, apare necesitatea unor
\emph{sisteme informatice integrate} capabile să gestioneze atât stocarea
distribuită a datelor, cât și rularea eficientă a modelelor matematice și
vizualizarea rezultatelor. Alegerea limbajului de programare și a arhitecturii
software are un impact direct și semnificativ asupra performanței, scalabilității
și mentenabilității unui astfel de sistem.

Lucrarea de față se înscrie în această direcție, propunând proiectarea și
realizarea unui sistem e-medicine pentru analizarea și monitorizarea datelor
epidemiologice, cu accent pe implementarea și analiza comparativă a modelului
SIR în patru medii de programare reprezentative pentru spectrul aplicațiilor
de calcul științific: \textbf{R}, \textbf{Julia}, \textbf{Python} și \textbf{SAS}.

\section{Obiectivele lucrării}
\label{sec:obiective}

Prezenta lucrare urmărește realizarea unui sistem coerent și cuprinzător, structurat
în jurul mai multor obiective principale:

\begin{enumerate}
    \item \textbf{Fundamentarea matematică riguroasă} a modelului SIR prin derivarea
    ecuațiilor din principii prime, demonstrarea proprietăților matematice fundamentale
    (existență, unicitate, pozitivitate, conservare) și analiza calitativă a
    soluțiilor (echilibre, stabilitate, numărul de reproducere de bază $\Rzero$).

    \item \textbf{Implementarea modelului SIR} în patru limbaje de programare ---
    R, Julia, Python și SAS --- utilizând metodele numerice adecvate pentru rezolvarea
    sistemelor de ecuații diferențiale ordinare (metoda Euler, Runge-Kutta de ordinul 4,
    metode adaptive).

    \item \textbf{Evaluarea comparativă a performanței} implementărilor pe multiple
    configurații hardware, urmărind metricile de timp de execuție, consum de memorie
    și scalabilitate în funcție de dimensiunea setului de date.

    \item \textbf{Proiectarea și implementarea infrastructurii distribuite} pentru
    stocarea datelor epidemiologice, bazată pe baza de date YugabyteDB rulată
    într-un cluster Kubernetes, asigurând disponibilitate ridicată și distribuție
    eficientă a datelor.

    \item \textbf{Dezvoltarea aplicației web e-medicine}, compusă dintr-un API
    backend implementat în Go și o interfață frontend Vue.js, care integrează
    toate componentele sistemului și oferă o perspectivă unificată asupra
    rezultatelor.
\end{enumerate}

\section{Contribuții originale}
\label{sec:contributii}

Contribuțiile originale ale acestei lucrări pot fi sintetizate după cum urmează:

\begin{itemize}
    \item \textbf{Analiză comparativă sistematică} a implementărilor modelului SIR
    în R, Julia, Python și SAS, acoperind atât corectitudinea numerică, cât și
    performanța computațională pe configurații hardware diverse.

    \item \textbf{Metodologie de benchmarking reproducibilă} pentru evaluarea
    performanței algoritmilor epidemiologici în limbaje de programare științifică,
    inclusiv protocoale de testare și metrici standardizate.

    \item \textbf{Arhitectură distribuită de stocare a datelor epidemiologice}
    bazată pe YugabyteDB și Kubernetes, cu schema relațională optimizată pentru
    interogări specifice domeniului epidemiologic.

    \item \textbf{Sistem web integrat} pentru vizualizarea și compararea
    rezultatelor modelării epidemiologice, combinând un API REST performant
    (Go) cu o interfață web modernă și reactivă (Vue.js).
\end{itemize}

\section{Structura lucrării}
\label{sec:structura}

Lucrarea este organizată în opt capitole, după cum urmează:

\textbf{Capitolul~\ref{ch:epidemiologie}} (\emph{Epidemiologie Computațională și
Modelarea Matematică}) prezintă fundamentele epidemiologiei matematice. Sunt
introduse modelele compartimentale (SI, SIS, SIR, SEIR, SIRS), cu accent pe
analiza matematică riguroasă a modelului SIR: derivare, teoreme de existență și
unicitate, analiza stabilității, numărul de reproducere de bază $\Rzero$ și
ecuația dimensiunii finale. Capitolul include și o trecere în revistă a metodelor
numerice utilizate pentru rezolvarea sistemelor de ecuații diferențiale ordinare.

\textbf{Capitolul~\ref{ch:tehnologii}} (\emph{Tehnologii și Instrumente Utilizate})
prezintă fundamentele teoretice și caracteristicile tehnice ale instrumentelor
utilizate: limbajele de programare R, Julia, Python și SAS, baza de date
distribuită YugabyteDB, sistemul de orchestrare Kubernetes, limbajul de
programare Go și framework-ul frontend Vue.js.

\textbf{Capitolul~\ref{ch:arhitectura}} (\emph{Arhitectura Sistemului E-Medicine})
descrie arhitectura globală a sistemului, cerințele funcționale și non-funcționale,
schema bazei de date distribuite, arhitectura API-ului backend și structura
componentelor frontend.

\textbf{Capitolul~\ref{ch:implementare}} (\emph{Implementarea Modelului SIR})
detaliază implementările practice ale modelului SIR în fiecare limbaj de
programare, inclusiv seturile de date utilizate, parametrii modelului și
interfața de integrare cu API-ul Go.

\textbf{Capitolul~\ref{ch:evaluare}} (\emph{Evaluarea Performanței și
Scalabilității}) prezintă metodologia de benchmarking, configurațiile hardware
utilizate și rezultatele comparative ale evaluărilor de performanță.

\textbf{Capitolul~\ref{ch:aplicatie}} (\emph{Aplicația Web E-Medicine})
detaliază implementarea aplicației web, documentând atât backend-ul Go, cât și
frontend-ul Vue.js, cu descrierea funcționalităților și a fluxurilor de utilizare.

\textbf{Capitolul~\ref{ch:concluzii}} (\emph{Concluzii și Direcții Viitoare})
sintetizează contribuțiile lucrării, formulează concluzii privind implementările
și infrastructura, și identifică direcții de dezvoltare ulterioară.

Lucrarea se încheie cu \textbf{Bibliografie} și \textbf{Anexe} care conțin
codul sursă complet al implementărilor, schema bazei de date, configurațiile
Kubernetes și documentația API-ului REST.
